\documentclass{article}
\usepackage[utf8]{inputenc}
\usepackage[portuguese]{babel}
\usepackage{amsmath}
\usepackage{amsfonts}
\usepackage{amssymb}
\usepackage{geometry}
\geometry{a4paper, margin=1in}

\title{MAC0344 Arquitetura de Computadores \\ Lista de Exercícios No. 3}
\author{Leonardo Heidi Almeida Murakami \\ NUSP: 11260186}
\date{\today}

\begin{document}
\maketitle

\section*{1. Identifique todas as dependências nas seguintes instruções}

Considere o seguinte código:
\begin{verbatim}
1: B = Y + 5
2: A = B + C
3: C = X * X
4: D = 0
5: E = D + 1
6: D = Y + 1
\end{verbatim}

As dependências identificadas são:

\begin{itemize}
    \item \textbf{Dependências Verdadeiras (RAW - Read After Write):}
    \begin{itemize}
        \item $1 \rightarrow^{v} 2$: A instrução 2 lê a variável $B$ após a instrução 1 escrever nela.
        \item $4 \rightarrow^{v} 5$: A instrução 5 lê a variável $D$ após a instrução 4 escrever nela.
    \end{itemize}
    \bigskip
    \item \textbf{Anti-dependências (WAR - Write After Read):}
    \begin{itemize}
        \item $2 \rightarrow^{\text{anti}} 3$: A instrução 3 escreve na variável $C$ após a instrução 2 ler o valor original de $C$.
        \item $5 \rightarrow^{\text{anti}} 6$: A instrução 6 escreve na variável $D$ após a instrução 5 ler o valor de $D$ definido pela instrução 4.
    \end{itemize}
    \bigskip
    \item \textbf{Dependências de Saída (WAW - Write After Write):}
    \begin{itemize}
        \item $4 \rightarrow^{\text{saída}} 6$: A instrução 6 escreve na variável $D$ após a instrução 4 também ter escrito em $D$.
    \end{itemize}
\end{itemize}

\newpage
\section*{2. Removendo as anti-dependências e dependências de saída}

As dependências falsas (anti-dependências e dependências de saída) podem ser removidas através de \textbf{renomeação de variáveis}, alocando novos registradores para os destinos conflitantes:

\begin{itemize}
    \item A anti-dependência $2 \rightarrow^{\text{anti}} 3$ ocorre porque $C$ é lida em 2 e escrita em 3. Renomeamos o destino da instrução 3 para $C1$.
    \item A dependência de saída $4 \rightarrow^{\text{saída}} 6$ e a anti-dependência $5 \rightarrow^{\text{anti}} 6$ ocorrem porque $D$ é escrita em 4 e 6, sendo lida em 5. Renomeamos o destino da instrução 6 para $D1$.
\end{itemize}

O código reescrito fica:
\begin{verbatim}
1: B = Y + 5
2: A = B + C
3: C1 = X * X
4: D = 0
5: E = D + 1
6: D1 = Y + 1
\end{verbatim}

Com essas renomeações, eliminamos todas as dependências falsas, mantendo apenas as dependências verdadeiras $1 \rightarrow^{v} 2$ e $4 \rightarrow^{v} 5$.

\end{document}